\documentclass[12pt]{article}
\usepackage[margin=1in]{geometry}
\usepackage{setspace}
\doublespacing

%% Bibliography
\usepackage{natbib}

%% Tables
\usepackage{booktabs, threeparttable, siunitx, caption}
\usepackage{tabu}
\usepackage{float}
\extrarowsep=0.1pt
\renewcommand{\arraystretch}{1.15}

%% Figures
\usepackage{graphicx}

%% Math & fonts
\usepackage{amsmath, amssymb}
\usepackage{mathpazo}

%% Cross-references
\usepackage{hyperref}
\hypersetup{colorlinks=true, linkcolor=black, citecolor=black, urlcolor=black}

%% Misc
\usepackage[dvipsnames]{xcolor}
\usepackage{microtype}

\def\sym#1{\ifmmode^{#1}\else\(^{#1}\)\fi}

% ─────────────────────────────────────────────────────────────────────────────
\begin{document}

% ── Title page ───────────────────────────────────────────────────────────────
\begin{titlepage}
\begin{center}
\vspace*{1.5cm}
{\large \textbf{Does Baseline Willingness-to-Pay Predict\\[0.3em]
Life Satisfaction Gains from Electrification?\\[0.3em]
Evidence from Rural Kenya}}\\[2em]
{\normalsize Extension of Lee, Miguel \& Wolfram (JPE 2019)}\\[1.5em]
{\normalsize \today}
\end{center}

\vspace{2em}
\begin{abstract}
\noindent
We ask whether a household's stated willingness-to-pay (WTP) for a grid
electricity connection at baseline predicts the treatment-on-the-treated
(TOT) gains it realises from connection. Using data from the randomised
experiment of Lee, Miguel \& Wolfram (2019), we merge a
contingent-valuation WTP measure from the 2014 baseline survey with
outcomes from the first follow-up round (2016). We estimate heterogeneous
TOT effects via IV2SLS---instrumenting actual connection with treatment
arm indicators and their interactions with a binary high-WTP
variable---on two outcomes: the number of appliances owned and a
normalised life satisfaction index. For appliances, low-WTP households
gain $0.211^{*}$ additional appliances and high-WTP households gain
$0.479^{**}$, but the differential $\hat{b}_3=0.268$ is not statistically
significant ($p=0.173$). For life satisfaction, both groups show small,
statistically insignificant TOT effects ($0.132$ and $0.144$
respectively), with a near-zero differential ($\hat{b}_3=0.012$,
$p=0.938$). These results suggest that prior expectations, as proxied by
stated WTP, do not meaningfully moderate the life satisfaction impact of
grid connection, while leaving open the possibility of heterogeneous
effects on material outcomes.
\end{abstract}
\end{titlepage}

% ── 1. Background and Motivation ─────────────────────────────────────────────
\section{Background and Motivation}
\label{sec:background}

A growing literature examines the welfare effects of household
electrification in developing countries. Lee, Miguel \& Wolfram (2019)
conduct a randomised experiment in rural Kenya in which households are
randomly offered grid connections at varying subsidy levels. Despite large
take-up responses to price subsidies, the authors find small or null
average effects on income, employment, and subjective wellbeing.

One mechanism that could generate heterogeneous wellbeing effects even
within an experiment with null average effects is an \emph{expectations
channel}: households that valued the connection more highly ex ante---as
revealed by higher stated WTP---may have had stronger prior expectations
of benefit. If electrification meets those expectations, high-WTP
households should exhibit larger life satisfaction gains than low-WTP
households; if it fails to meet them, the reverse could hold.

This note tests the expectations channel directly by interacting actual
grid connection (instrumented by the treatment arms) with a baseline WTP
measure. The specification follows Lee, Miguel \& Wolfram (2019, JEP)
Appendix Tables A3--A4, which interact grid connection with a high-SES
composite index. We substitute the SES index with a binary high-WTP
indicator derived from the contingent-valuation exercise conducted at
baseline, and extend the analysis from wellbeing alone to also cover
appliance ownership.

% ── 2. Data and Specification ────────────────────────────────────────────────
\section{Data and Specification}
\label{sec:data}

\subsection*{Data}

We use three datasets from Lee, Miguel \& Wolfram (2019). The baseline
survey (\texttt{demand.dta}, 2014) contains the WTP elicitation and
treatment assignment. The first follow-up survey (\texttt{impacts\_1.dta},
2016, Round~1) contains the outcome variables used in this analysis. Both
datasets share a unique household identifier (\texttt{reppid}), enabling
a 1:1 merge. We drop households already connected at baseline
(\texttt{base\_connected}$=1$) and restrict to adult respondents
(\texttt{child}$=0$), following the original papers. The resulting sample
contains $N=2{,}217$ observations.

\subsection*{WTP Measure}

At baseline, each household was randomly assigned a hypothetical connection
price (\texttt{WTP\_amt} $\in \{0,\,10{,}000,\,\ldots,\,75{,}000\}$~KES)
and asked whether it would accept (\texttt{WTP\_r1}$\in\{0,1\}$). We
define stated WTP as:
\begin{equation}
\label{eq:wtp}
\text{WTP-stated}_i \;=\; \texttt{WTP\_amt}_i
  \;\times\; \mathbf{1}[\texttt{WTP\_r1}_i = 1].
\end{equation}
Since the elicitation price is randomly assigned, stated WTP has an
exogenous component. The sample median of WTP-stated equals zero (56.4\,\%
of households declined the offer), so
$\texttt{wtp\_high}_i = 1$ effectively identifies the 43.6\,\% of
households that accepted at any positive price.

\subsection*{Outcomes}

Our primary outcomes are: (1)~\texttt{number\_appliances}, the count of
appliances owned by the household, which serves as a measure of material
living standards; and (2)~\texttt{life\_index}, a normalised life
satisfaction index. Both outcomes are available in Round~1 (2016). The
life satisfaction index is standardised to mean zero and unit standard
deviation in the control group.

\subsection*{Specification}

We estimate treatment-on-the-treated (TOT) effects via two-stage least
squares (IV2SLS), following \texttt{tab\_A3\_A4.do} in the replication
archive of Lee, Miguel \& Wolfram (2019, JEP). The second-stage equation
is:
\begin{equation}
\label{eq:spec}
y_{ic} \;=\; \alpha
  \;+\; b_1\,\text{Connected}_{ic}
  \;+\; b_2\,\text{HighWTP}_{i}
  \;+\; b_3\,(\text{Connected}_{ic} \times \text{HighWTP}_{i})
  \;+\; \boldsymbol{\delta}'\mathbf{X}_{ic}
  \;+\; \varepsilon_{ic},
\end{equation}
where $\text{Connected}_{ic}$ is actual grid connection status and
$\text{HighWTP}_i = \mathbf{1}[\text{WTP-stated}_i > \text{median}]$.
The instruments are the treatment arm indicators
(\texttt{treat\_low}, \texttt{treat\_med}, \texttt{treat\_full}) and
their interactions with \texttt{HighWTP}, mirroring the excluded
instruments in \texttt{tab\_A3\_A4.do}. $\mathbf{X}_{ic}$ includes
community controls (county, market proximity, transformer timing, baseline
connection rate, population) and household controls (gender, age, housing
quality, baseline energy spending). Column~(2) of Table~\ref{tab:main}
reports $\hat{b}_1$ (the TOT for low-WTP households); column~(3) reports
$\hat{b}_1+\hat{b}_3$ (the TOT for high-WTP households); column~(4)
reports $\hat{b}_3$ (the differential). Standard errors are clustered at
the community level (\texttt{siteno}).

% ── 3. Results ───────────────────────────────────────────────────────────────
\section{Results}
\label{sec:results}

Figure~\ref{fig:coefplot} plots the TOT coefficients for low-WTP (hollow
circles) and high-WTP (filled circles) households, with 95\,\% confidence
intervals, for both outcomes. For appliances (Panel~A), there is a
visible positive gradient: the high-WTP point estimate ($0.479$) is
approximately twice the low-WTP estimate ($0.211$), though the
differential is not statistically distinguishable from zero given the
wide confidence intervals. For life satisfaction (Panel~B), the two
groups' estimates are nearly identical and straddle zero.

\begin{figure}[H]
  \centering
  \includegraphics[width=0.92\textwidth]{%
    ../figures/wellbeing_wtp_heterogeneity_2.pdf}
  \caption{TOT coefficients on outcomes by baseline WTP group (IV2SLS).
    Hollow circles: low-WTP households (WTP-stated $\leq$ median).
    Filled circles: high-WTP households (WTP-stated $>$ median).
    Horizontal bars are 95\,\% confidence intervals. Vertical dashed
    line at zero. Sample: Round~1 (2016), $N=2{,}217$.
    See Section~\ref{sec:data} for specification.}
  \label{fig:coefplot}
\end{figure}

Table~\ref{tab:main} presents the full regression results. For appliances
(Panel~A, row~A1), the low-WTP TOT is $0.211^{*}$ (SE$=0.111$,
$p=0.057$), the high-WTP TOT is $0.479^{**}$ (SE$=0.207$, $p=0.021$),
and the differential $\hat{b}_3=0.268$ (SE$=0.197$, $p=0.173$). The
point estimates suggest that connected high-WTP households accumulate
roughly one-quarter of an additional appliance more than low-WTP
households, but this difference is not statistically significant at
conventional levels.

For life satisfaction (Panel~B, row~B1), the low-WTP TOT is $0.132$
(SE$=0.099$, $p=0.181$) and the high-WTP TOT is $0.144$ (SE$=0.106$,
$p=0.174$). The differential $\hat{b}_3=0.012$ (SE$=0.160$, $p=0.938$)
is negligible in magnitude and far from statistically significant. The
near-zero differential is consistent with the null average wellbeing
effects reported in the original paper, and suggests that the expectations
channel does not operate through life satisfaction.

\clearpage
\newgeometry{left=0.9in, right=0.9in, top=1in, bottom=1in}
\begin{singlespace}

% ── Table 1 (inlined from tab_wellbeing_main_2.tex) ──────────────────────────
\begin{table}[H]
\centering
\caption{Heterogeneous TOT effects of grid connection,
  by baseline willingness-to-pay}
\label{tab:main}
\vspace{-0.5em}
\begin{tabu} to \linewidth {X[20,l] X[5,c] X[5,c] X[5,c] X[5,c] X[4,c]}
\toprule
 & Control & Low-WTP & High-WTP & TOT\,$\times$ & \\[-0.25em]
 & [SD] & TOT & TOT & HighWTP & $N$ \\[-0.25em]
 & (1) & (2) & (3) & (4) & (5) \\
\midrule \\[-1.2em]

%% ── Panel A ──────────────────────────────────────────────────────────────
\multicolumn{6}{l}{\emph{Panel A: Appliance Outcomes}} \\[0.5em]

\hspace{2.5mm} A1. Number of appliances owned
  & 1.832 & $0.211$\sym{*} & $0.479$\sym{**} & $0.268$ & 2217 \\
& [1.340] & $(0.111)$ & $(0.207)$ & $(0.197)$ & \\[0.8em]

%% ── Panel B ──────────────────────────────────────────────────────────────
\multicolumn{6}{l}{\emph{Panel B: Wellbeing Outcomes}} \\[0.5em]

\hspace{2.5mm} B1. Normalised life satisfaction
  & -0.000 & $0.132$ & $0.144$ & $0.012$ & 2217 \\
& [1.000] & $(0.098)$ & $(0.106)$ & $(0.160)$ & \\[0.8em]

\bottomrule
\end{tabu}

\begin{minipage}{\linewidth}
\footnotesize
\setlength{\parskip}{2pt}
\textit{Notes}: Column~(1) reports the control-group mean with the standard
deviation in brackets. Columns~(2) and~(3) report TOT coefficients estimated
by IV2SLS. Instruments: \texttt{treat\_low}, \texttt{treat\_med},
\texttt{treat\_full} and their interactions with \texttt{HighWTP}.
Column~(2) is $\hat{b}_1$ (low-WTP TOT) and column~(3) is
$\hat{b}_1 + \hat{b}_3$ (high-WTP TOT); standard errors computed by the
delta method. Column~(4) reports $\hat{b}_3$ (differential effect) with its
standard error in parentheses. Sample: adults not pre-connected at baseline
(R1, 2016). \text{HighWTP} $= 1$ if stated WTP $>$ sample median ($43.6\,\%$
of sample). Controls: county, market proximity, transformer timing, baseline
connection rate, population, gender, age, housing quality, baseline energy
spending. Standard errors clustered at community level (\texttt{siteno}).
Asterisks indicate significance (two-tailed):
$^{*}p<0.10$;\;$^{**}p<0.05$;\;$^{***}p<0.01$.
\end{minipage}
\end{table}
% ─────────────────────────────────────────────────────────────────────────────

\end{singlespace}
\restoregeometry

% ── 4. Caveats ───────────────────────────────────────────────────────────────
\section{Caveats}
\label{sec:caveats}

\paragraph{IV exclusion restriction.}
The IV2SLS estimates are valid only if the treatment arm indicators affect
outcomes solely through actual grid connection (the exclusion restriction).
This would be violated if, for example, being assigned to a subsidised
connection arm changed household behaviour or expectations independently
of whether connection actually occurred. In the context of this experiment,
the exclusion restriction is plausible---treatment arms were randomly
assigned and the primary mechanism through which assignment would affect
appliance ownership or life satisfaction is through the connection
itself---but cannot be formally tested.

\paragraph{WTP measure is binary and partially random.}
Because each household sees a single randomly assigned elicitation
price, we cannot recover a complete demand curve at the individual
level. The \texttt{wtp\_high} indicator conflates the household's
underlying valuation with the random price draw: a household assigned
a high price that accepts may differ systematically from one assigned a
low price that accepts. The median split on WTP-stated is therefore a
coarse proxy for true willingness-to-pay.

\paragraph{R1-only sample.}
The \texttt{\_2} analysis is restricted to Round~1 (2016, $N=2{,}217$),
following the availability of the \texttt{number\_appliances} variable.
This limits the pooled sample size relative to the original ITT analysis
and precludes panel-based specifications or comparison with Round~2
outcomes. Results may not generalise to long-run effects measured in R2
(2017).

\paragraph{Null average effects constrain power.}
The original papers find small average effects on \texttt{life\_index}.
With such small signal, detecting heterogeneity of plausible magnitude
would require a substantially larger sample or stronger variation in WTP.
The appliance outcome shows larger point estimates, but the differential
remains insignificant due to wide confidence intervals from the IV
first-stage noise.

% ── References ───────────────────────────────────────────────────────────────
\clearpage
\begin{singlespace}
\begin{thebibliography}{9}

\bibitem[Benjamini and Hochberg(1995)]{benjamini1995}
Benjamini, Yoav, and Yosef Hochberg. 1995.
``Controlling the False Discovery Rate: A Practical and Powerful
Approach to Multiple Testing.''
\textit{Journal of the Royal Statistical Society, Series~B}
57~(1): 289--300.

\bibitem[Lee, Miguel, and Wolfram(2019)]{lee2019jpe}
Lee, Kenneth, Edward Miguel, and Catherine Wolfram. 2019.
``Experimental Evidence on the Economics of Rural Electrification.''
\textit{Journal of Political Economy} 128~(4): 1523--1565.

\bibitem[Lee, Miguel, and Wolfram(2020)]{lee2019jep}
Lee, Kenneth, Edward Miguel, and Catherine Wolfram. 2020.
``Does Household Electrification Supercharge Economic Development?''
\textit{Journal of Economic Perspectives} 34~(1): 122--144.

\end{thebibliography}
\end{singlespace}

\end{document}
